\changecolours{ScoutGreen}
\section[Revisão e Objetivos]{Ponte Pedagógica e Objetivos de Aprendizagem}

% ------------------------------------------------------------------------------
\twocolframe[title={Retrospectiva da Aula 05 e o Novo Desafio},leftcol=7.1cm,rightcol=6.9cm]{%
  \alert{O Que Consolidamos na Aula 05?}\par
  \itemseps[topsepi=0.5mm,itemsepi=1mm,labelsep=2mm,leftmargini=3mm]
  \begin{itemize}
    \item \textbf{Relacionamentos no MER:} Associações semânticas como subconjunto do produto cartesiano ($R \subseteq E_1 \times E_2$).
    \item \textbf{Cardinalidade:} Razões $1:1$, $1:N$ e $N:M$ derivadas estritamente das regras de negócio.
    \item \textbf{Participação:} Linha dupla (total/obrigatória) e simples (parcial/opcional).
    \item \textbf{Atributos Descritivos:} Alocados no losango de ligação.
  \end{itemize}
}{%
  \alert{O Desafio Central da Aula 06}\par
  \itemseps[topsepi=0.5mm,itemsepi=1mm,labelsep=2mm,leftmargini=3mm]
  \begin{itemize}
    \item \textbf{Entidades Dependentes:} E quando uma entidade do mundo real \textbf{não possui atributos suficientes} para formar uma chave primária própria e autônoma?
      \begin{itemize}\setlength{\itemsep}{0.3mm}
        \item Ex: As portas de rede de um roteador.
        \item Ex: Os dependentes de um funcionário.
      \end{itemize}
    \item \textbf{Taxonomias e Herança:} Como representar hierarquias do tipo ``é um'' sem duplicar atributos comuns?
  \end{itemize}
}

% ------------------------------------------------------------------------------
\twocolframe[title={Objetivos de Aprendizagem da Aula 06},leftcol=7.1cm,rightcol=6.9cm]{%
  \alert{Competências Teóricas e Conceituais}\par
  \itemseps[topsepi=0.5mm,itemsepi=1mm,labelsep=2mm,leftmargini=3mm]
  \begin{itemize}
    \item \textbf{Entidades Fracas:} Dominar os conceitos de dependência existencial e de identificação.
    \item \textbf{Chave Parcial (Discriminador):} Formar a chave primária composta com a entidade proprietária.
    \item \textbf{Notação de Peter Chen:} Aplicar retângulos duplos, losangos duplos e sublinhado tracejado.
    \item \textbf{Modelo EER:} Projetar hierarquias com superclasses, subclasses e herança de atributos.
  \end{itemize}
}{%
  \alert{Competências Metodológicas e Práticas}\par
  \itemseps[topsepi=0.5mm,itemsepi=1mm,labelsep=2mm,leftmargini=3mm]
  \begin{itemize}
    \item \textbf{Restrições Estruturais EER:} Classificar especializações em Disjuntas ($d$) ou Sobrepostas ($o$), Totais ou Parciais.
    \item \textbf{Modelagem em Telemática:} Aplicar conceitos a equipamentos de rede, interfaces e topologias.
    \item \textbf{Laboratório Computacional:} Validar regras em memória via metamodelo em Python orientado a objetos.
    \item \textbf{Alinhamento AVP1:} Esboçar entidades fracas e subclasses para o projeto semestral.
  \end{itemize}
}
